\chapter{Reproducing the measurements}
\label{app:reproducibility}

This appendix describes the steps for reproducing the measurements of this
work. The procedure covers installing the tools, running the code, and
generating the plots and the report. The commands concern the system of
section~\ref{sec:meth-machine} and are run from the root of the repository
(\url{https://github.com/ioannisam/topk}).

\section{Software requirements}
\label{sec:repro-requirements}

All three implementations require a C++ compiler that supports the C++20
standard, as well as cmake version 3.18 or newer. Vector instructions are not a
compilation requirement. The choice between AVX-512, AVX2 and scalar execution
is made at run time, based on the capabilities reported by the processor.
Because the measurements were made on a machine with AVX-512, its absence leads
to a different code path and the results are not directly comparable.

The graphics unit requires the CUDA toolkit, which includes the nvcc compiler
and the CUB library. Power measurements additionally require the driver's NVML
library.

The neural processing unit requires the XRT runtime stack and the XDNA plugin.
Compiling the kernels needs MLIR-AIE in a Python environment. This environment
must reside at the root of the repository under the name npu\_venv, since that
is where the run scripts look for it. The user must belong to the device's
access groups and have a sufficient locked-memory limit; otherwise the device
is not recognised. The XILINX\_XRT variable must point to the XRT installation.

Energy measurement reads the RAPL counters from /sys/class/powercap, which
requires administrator privileges. In a single run without them, the
measurement is explicitly declared unavailable instead of recording zeros
(section~\ref{sec:meth-energy}). The automated sweep, in contrast, is aborted
entirely to avoid producing an empty dataset. Finally, the analysis
infrastructure requires the repository's Python libraries. The corresponding
environment is created automatically on the first run.

\section{Compilation}
\label{sec:repro-build}

All implementations are compiled with a single command:

\begin{verbatim}
make build-all
\end{verbatim}

Each architecture can also be compiled on its own, with the targets build-cpu,
build-gpu and build-npu, which is useful on a system that does not have all
three units. The resulting executables are placed in the build directory, two
per architecture: one for the top-k measurements and one for the ceiling
measurements of section~\ref{sec:meth-roofline}.

Compilation for the neural processing unit additionally produces the hardware
configuration files. Because their absence leads to the corresponding
measurements being silently skipped, their presence is checked explicitly:

\begin{verbatim}
make verify-artifacts
\end{verbatim}

The check examines the two configuration files and the three top-k executables.
It presupposes a full compilation, so it fails automatically if only one
architecture has been compiled.

\section{Running the experiments}
\label{sec:repro-run}

The recommended way to reproduce the results is to run the whole procedure with
a single command:

\begin{verbatim}
make benchmark GPU_W=50
\end{verbatim}

This command ensures the correct sequence of steps. If the conditions are
pinned after the fact or old executables are used, the results lose their
comparability. The command enforces the proper order and prevents omissions.

The sequence comprises six stages. First, the old results and executables are
deleted. Compilation and the check of the files follow
(section~\ref{sec:repro-build}). In the third stage the system conditions are
pinned and access to the energy counters is granted. The four measurement
sweeps follow, then the generation of the plots and, finally, the restoration
of the system.

The procedure first checks the availability of sudo and of the npu\_venv
environment. If either is missing, the command terminates immediately, without
altering the system or the data.

Administrator privileges are requested once and remain active. The conditions
are always restored at the end, even in case of an error or an interruption by
the user, preventing the system from being left in a pinned state.

The GPU\_W argument sets the power limit of the graphics unit, which must
remain constant (section~\ref{sec:meth-conditions}). The procedure deletes the
existing results immediately, so the user must first keep copies of any data
they want to preserve.

\subsection{Running individual steps}
\label{sec:repro-steps}

For a partial or repeated run, each stage is also available separately. The
system conditions are controlled with:

\begin{verbatim}
make pin ARGS=50      # governor, boost off, GPU cap 50 W
make pin-show         # report the current state
make unpin            # restore governor, boost, GPU cap
\end{verbatim}

The four independent sweeps write their raw results to separate files:

\begin{verbatim}
make measure-cases      # timing over the parameter space
make measure-energy     # energy, requires RAPL access
make measure-dists      # input-distribution sensitivity
make measure-roofline   # bandwidth, compare-exchange roofs
\end{verbatim}

The targets accept additional arguments through the ARGS variable, restricting
the parameter space. This makes it easy to check the installation quickly. In a
manual run, the user takes responsibility for the correct order and for keeping
the conditions constant.

\subsection{Execution parameters}
\label{sec:repro-cli}

Table~\ref{tab:cli} summarises the parameters of each executable, the permitted
values, the defaults and the values used in the sweeps.

The defaults often differ from the final sweep values. The value range was set
up to $10^9$ instead of the default, to avoid an excessive accumulation of ties
in inputs of millions of elements. Many ties would artificially favour the
algorithms that reject elements quickly.

This range applies to four of the five types. For half, the executable clamps
the bounds to $\pm 65504$, the largest finite value of the type. In this case
ties remain unavoidable, because of the narrow range and the inability of half
to represent all integers above $2048$. This property stems from the type
itself and directly affects the interpretation of the results.

The random seed is generated deterministically from the q and k of each case.
This ensures reproducibility without the same input ever being given under
different conditions.

The main sweep uses the values of the fourth column. The special sweeps
restrict the parameters: the distribution sweep examines q from 20 to 24, the
float and half types, and the k extremes 8 and 131072, with two seeds per case.
The energy sweep examines q from 16 to 24 with three repetitions, skipping the
correctness check so as not to inflate the consumed power needlessly.

\begin{table}[H]
    \centering
    \footnotesize
    \setlength{\tabcolsep}{4pt}
    \begin{tabularx}{\textwidth}{@{} l >{\raggedright\arraybackslash}X l >{\raggedright\arraybackslash}X >{\raggedright\arraybackslash}X @{}}
        \toprule
        \textbf{Parameter} & \textbf{Values} & \textbf{Default} & \textbf{In the sweeps} & \textbf{Meaning} \\
        \midrule
        q & integer 0--63 & --- & 3--24 & Input size $n = 2^q$ (only mandatory) \\
        k & positive integer & n & 8, 256, 4096, 131072 & Number of requested elements (up to n) \\
        mode & min, max & max & max & Select smallest or largest elements \\
        dtype & int, uint, float, double, half & int & all five & Element type \\
        algo & bitonic, map\_reduce, gt & bitonic & all three & Algorithm to run \\
        run & full, trunc, both & trunc & trunc & Variant of bitonic network \\
        seed & $\ge 0$ integer & 42 & from q, k & Random-number generator seed \\
        dist & uniform, normal, trimodal, sorted, reverse, adversarial & uniform & uniform (all 5 in sweep) & Input distribution \\
        min/max & integers & 0/1000 & 0/$10^9$ & Value range (max $\ge$ min) \\
        threads & positive integer & 0 (auto) & 8 & Number of threads (ignored by NPU) \\
        verify & true, false & false & true & Check against verification \\
        debug & true, false & true & false & Verbose diagnostic output \\
        \bottomrule
    \end{tabularx}
    \caption{The execution parameters and their defaults.}
    \label{tab:cli}
\end{table}

\section{Generating the plots}
\label{sec:repro-plots}

The analysis infrastructure reads the raw results, converts them into CSV
tables and generates the plots:

\begin{verbatim}
make plot ARGS="--plot all"
\end{verbatim}

The -{-}plot argument selects the plots and -{-}input the sweep (the default is
the timing sweep). Because the other plots use their own files, the combined
targets are recommended (measurement and generation):

\begin{verbatim}
make benchmark-cases
make benchmark-energy
make benchmark-dists
make benchmark-roofline
\end{verbatim}

The files are organised by data type and category. Plots that concern only the
hardware (e.g.\ the bandwidth scale) are stored in a separate directory, since
they do not depend on the type.

The -{-}energy-metric parameter selects between gross and net energy. The
sweeps of this work use the net figure (section~\ref{sec:meth-energy}). The
-{-}error-bars parameter sets the dispersion band (standard deviation, deciles
or none). In the automated run, the bands are disabled entirely.
